\section{Validation, reproducibility, and remaining proof obligations}
\label{sec:validation}

\subsection{What was checked}

The delivered production verifier constructs every support from the stated
ordered plans and matrix weights. Rigorous integer branch bounds cover
$173{,}846{,}817$ Cartesian support tuples, directly evaluate $7{,}560$
tuple norms, and solve $18$ nonuniform hull problems with exact rational
arithmetic. Every pruned block is bounded above, counted, and included in
an exact sum-of-squared-norms checksum. The code validates the projection
inequalities for the sharp witnesses. It does not sample a continuous set
of entropy vectors or use floating-point optimization as a proof.

The independent implementation reconstructs the supports and evaluates
all tuple norms without the production branch pruning. It obtains the
same tuple counts, maxima, checksums and certified rational bounds.
Its support ordering may differ, so witness indices need not be identical.
The twelve-group certificate for the newer backend is an upper bound;
only the other three certificates are asserted to be sharp for their menus.

The schedule verifier uses convergent rational series for logarithms and
$\pi$, and integer square-root enclosures. It checks every prefix gate,
the infinite-tail gates, the complete paid time recurrence, storage, and
strict separation from the rounded eventual constants. No binary
floating-point result is used to accept a schedule. The independent
schedule check also verifies the classical packing gates by integer-power
comparisons and recomputes its constants at high precision.

A small exact evaluator checks all four plans on sixteen finite mode-graph
fixtures over $(\mathbb Z/64\mathbb Z)[V]/(V^2)$, including a zero divisor
and a nonzero nilpotent. Its $64$ plan evaluations agree with literal
enumeration over $1{,}024$ assignments. It records $495$ actual child-entry
calls while constructing the caches, and checks two invalid-input controls.
This tests the contraction and selection identities over a nonfield; it is
not an implementation of the asymptotically fast matrix backend.

The modular-rank compiler passes $945$ finite-field factorization cases,
nine crossing-identity cases, four invalid-input rejections, an
exceptional-prime rank check, $2{,}049$ nonnegative CRT cases,
$4{,}097$ signed CRT cases, and $50$ signed-network cases followed by
two-adic reduction. One negative control demonstrates why the least
nonnegative CRT representative cannot simply be reduced modulo $2^\nu$.
Another demonstrates that exponentially many small-rank pieces violate
the uniform polynomial-piece interface.

All delivered production drivers were run both normally and with
\texttt{python3 -O}; their mathematical results agree. These checks use
explicit exceptions rather than removable assertions. The full replay,
artifact hashes and dependency boundaries are recorded in
\path{results/} and \path{PROOF_STATUS.json}.

\subsection{How to reproduce the finite claims}

From the archive's top-level directory, run
\begin{lstlisting}
python3 run_checks.py
python3 run_checks.py --optimized
python3 audit/independent/replay_menu_certificate.py --help
python3 build_report.py
python3 verify_bundle.py
\end{lstlisting}
The core certificate, schedule, comparison, algebra and modular checks use
Python's standard library. The optional independent exhaustive replay
uses NumPy; its independent high-precision schedule check uses mpmath.
The PDF build requires a standard TeX installation with pdfLaTeX and the
packages listed in \path{report.tex}. The exact command lines and expected
outputs are in \path{README.txt}.

\subsection{What has not been established}

The new finite minimax theorem, menu certificates, matrix-coordinate
transfer argument and schedule analysis are supplied here. The exact
deformation, Hermite recovery, Cauchy mode identity, entry interfaces and
SAT reductions remain imported from the cited sources. This is therefore
a relative improvement of those constructions, with those premises made
explicit, rather than an independent proof of their entire SAT theorem.

We did not build a complete SAT solver, benchmark a SAT instance, implement
the large fixed fast-matrix schemes, or independently rebuild the supplied
Lean formalization. No claim of practical crossover, global optimality
over all contraction plans, polynomial-space improvement, or a smaller
worst-case logarithmic order is made. The actual-rank criterion in
Section~\ref{sec:order-research} is an open theorem obligation, not a
consequence of the small finite probes.

The archive is suitable for review and integration as a new research
packet. Its exact finite certificates and explicit assumptions make the
proposed improvement concrete; any future claim of a better logarithmic
power should discharge the rank criterion or supply a different paid
uniform interface.
